\chapter*{บทที่ 7 สรุป ปัญหา และข้อเสนอแนะ}\addcontentsline{toc}{chapter}{บทที่ 7 สรุป ปัญหา และข้อเสนอแนะ}

\section*{7.1 สรุปผล}\addcontentsline{toc}{section}{7.1 สรุปผล}

คณะผู้จัดทำติดตั้ง WordPress บนเซิร์ฟเวอร์คลาวด์ด้วยตนเองครบทั้ง 3 ขั้นตามโจทย์ เว็บไซต์เข้าถึงได้ที่ \url{https://elephotelephus.site} ผ่าน HTTPS, SSH ใช้กุญแจเท่านั้นและปิด root, geo-firewall รับเฉพาะ IP ไทย 1,167 ช่วงและบล็อกการทดสอบจากต่างประเทศได้ทั้งหมด และเว็บไซต์ประชาสัมพันธ์สาขาใช้งานได้จริงพร้อมข่าวจากคณะที่อัปเดตเอง

ส่วนเพิ่มเติมคือระบบ secdash ที่ทำให้เห็นการโจมตีที่เกิดขึ้นจริง เช่น บอทสแกนไฟล์ลับและการไล่หาชื่อผู้ใช้ และใช้ AI สรุปเป็นภาษาไทยทุกชั่วโมง บทเรียนสำคัญคือการป้องกันหลายชั้นทำให้ชั้นหนึ่งที่พลาดไม่กลายเป็นความเสียหาย และคำแนะนำจาก AI ต้องตรวจกับระบบจริงก่อนนำไปใช้

\section*{7.2 ปัญหาและแนวทางแก้ไข}\addcontentsline{toc}{section}{7.2 ปัญหาและแนวทางแก้ไข}

ตารางที่ 7.1 ปัญหาที่พบระหว่างทำโครงงาน

\begin{longtable}[]{@{}>{\raggedright\arraybackslash}p{0.242\linewidth - 2\tabcolsep}>{\raggedright\arraybackslash}p{0.364\linewidth - 2\tabcolsep}>{\raggedright\arraybackslash}p{0.364\linewidth - 2\tabcolsep}@{}}
\toprule\noalign{}
ปัญหา & สาเหตุ & วิธีแก้ \\
\midrule\noalign{}
\endhead
\bottomrule\noalign{}
\endlastfoot
setup.sh หยุดทำงานเงียบๆ หลังใส่ Zone ID & ฟังก์ชันสุ่มรหัสผ่านใช้ pipe ที่เกิด SIGPIPE เมื่อเปิด set -e & อ่านจาก /dev/urandom เป็นก้อนขนาดคงที่แทน และตรวจ token กับ Cloudflare ก่อนทำต่อ \\
ตั้งค่าภาษา WordPress ไม่สำเร็จ & ไม่พบไฟล์ภาษา th ในขั้นติดตั้ง & ติดตั้งเป็น en\_US แล้วเปลี่ยนภาษาภายหลังในหน้าตั้งค่า \\
wp-cli สร้าง rewrite rule ไม่ได้ & สิทธิ์ไดเรกทอรีทำงานไม่ตรงกับผู้ใช้ www-data & เรียก wp-cli เป็น www-data ในโฟลเดอร์เว็บ และไม่ใช้ -\/-hard \\
เว็บตอบ NXDOMAIN & มี A record เฉพาะ cs ยังไม่มี @ & เพิ่ม A record ของ @ และล้าง DNS cache \\
แหล่ง IP ไทย (ipdeny) ส่งรายการว่าง & ดาวน์โหลดไม่สำเร็จ สคริปต์จึงไม่ยอมเปิด firewall & เปลี่ยนไปใช้ไฟล์สถิติของ APNIC และเก็บ ipdeny เป็นแหล่งสำรอง \\
สคริปต์หา IP ของผู้ดูแลไม่เจอ & sudo ลบตัวแปร SSH\_CLIENT & ใช้ค่าสำรองจาก who -m \\
ขั้น 2FA ถูกนับเป็น login ผิด & หน้า 2FA ตอบ 200 เหมือนรหัสผิด & ตัดออกเมื่อ 2FA สำเร็จตามมาใน 5 นาที \\
IP ของทีมถูกนับเป็นผู้โจมตี & ตัวแก้ไขบทความเรียก API รายชื่อผู้ใช้ & IP ที่ login สำเร็จถือเป็นคนในทีม \\
เรียก AI ไม่สำเร็จ (404, 503, timeout) & โมเดลเดิมถูกปลด เซิร์ฟเวอร์ AI ไม่ว่าง และรุ่นใหม่ตอบช้า & เลือกรุ่นอัตโนมัติ รอนานขึ้น และสลับรุ่นเมื่อมีปัญหา \\
DBeaver ต่อฐานข้อมูลไม่ได้ & ตั้งให้ต่อตรงพอร์ต 3306 ซึ่งไม่ได้เปิด & ใช้ SSH tunnel และตั้ง allowPublicKeyRetrieval \\
\end{longtable}

\section*{7.3 ข้อจำกัด}\addcontentsline{toc}{section}{7.3 ข้อจำกัด}

\begin{itemize}
\tightlist
\item
  geo-firewall กันผู้โจมตีที่ใช้ IP ไทย (VPN ไทย หรือเครื่องในไทย) ไม่ได้ จึงต้องพึ่งชั้นอื่นด้วย
\item
  คนไทยที่อยู่ต่างประเทศและบริการตรวจเว็บจากต่างประเทศ (เช่น SSL Labs) เข้าเว็บไม่ได้
\item
  การจัดประเภทการโจมตีใช้รูปแบบที่กำหนดไว้ การโจมตีแบบใหม่อาจไม่ถูกตรวจพบ
\item
  เซิร์ฟเวอร์มีเครื่องเดียว ถ้าเครื่องล่มเว็บจะหยุดให้บริการ
\end{itemize}

\section*{7.4 ข้อเสนอแนะ}\addcontentsline{toc}{section}{7.4 ข้อเสนอแนะ}

\begin{enumerate}
\def\labelenumi{\arabic{enumi}.}
\tightlist
\item
  ติดตามแพตช์ถาวรของ CVE-2026-115560 แล้วอัปเดต PHP ทันทีที่ออก และรัน security-check.sh เป็นประจำ
\item
  ใช้ reverse proxy หรือ cache ของ Cloudflare เพื่อเพิ่มความเร็วและซ่อน IP จริง โดยย้ายการกรองประเทศไปไว้ที่ Cloudflare
\item
  ส่งแจ้งเตือนทันที (เช่น อีเมลหรือ LINE) เมื่อระดับความเสี่ยงเป็น ``สูง''
\item
  เก็บไฟล์สำรองไว้นอกเซิร์ฟเวอร์ เช่น Cloud Storage กันกรณีเครื่องเสียหายทั้งเครื่อง
\item
  เพิ่ม Web Application Firewall เช่น ModSecurity กับชุดกฎ OWASP CRS
\end{enumerate}
